% ============================================================
%  Book chapter 18: Binary Data
% ============================================================
\chapter{Binary Data}

\textbf{What this chapter gives you:} the tools to read and write raw bytes with fixed layouts
(\code{struct}), to show bytes as hex or base64 (\code{binascii}, \code{base64}), and to work on
single bits with Python's bit operators. You need this for model weight files, embedding dumps,
image bytes in API payloads and network packets.

\tikzset{c18b/.style={draw=#1,fill=#1!12,minimum width=8mm,minimum height=7mm,inner sep=0pt,
    font=\ttfamily\small},
  c18bit/.style={draw=cgray,minimum width=6mm,minimum height=6mm,inner sep=0pt,font=\ttfamily\small}}

% ------------------------------------------------------------
\section{Convert Binary Data with struct}

Python has many text tools and few binary ones. The standard \code{struct} module is the main
one. It converts between Python values and a \code{bytes} object laid out like a C \code{struct}:
fixed-size fields, one after another, no separators, no labels.

\begin{center}
\begin{tikzpicture}[box/.style={draw=#1,fill=#1!10,rounded corners=3pt,minimum height=9mm,
    font=\ttfamily\small,inner sep=4pt},>=Stealth]
  \node[box=cgreen,align=center] (v) at (0,0) {Python values\\[-1pt]{\scriptsize\sffamily 1, 4, 3}};
  \node[box=cblue] (b) at (6,0) {bytes};
  \node[box=cgreen,align=center] (t) at (11,0) {tuple\\[-1pt]{\scriptsize\sffamily (1, 4, 3)}};
  \draw[->,thick] (v.east) -- node[above,font=\ttfamily\footnotesize]{struct.pack(fmt, ...)} (b.west);
  \draw[->,thick] (b.east) -- node[above,font=\ttfamily\footnotesize]{struct.unpack(fmt, b)} (t.west);
  \node[font=\sffamily\scriptsize,text=cgray] at (6,-0.9) {\code{fmt} must be the same on both sides};
\end{tikzpicture}
\end{center}
% MD-BEGIN
% ```mermaid
% flowchart LR
%     V["Python values (1, 4, 3)"] -- "struct.pack(fmt, ...)" --> B["bytes"]
%     B -- "struct.unpack(fmt, b)" --> T["tuple (1, 4, 3)"]
% ```
% *`fmt` must be the same on both sides.*
% MD-END

Suppose a small file starts with a 12-byte header: a 4-byte tag, a 2-byte version, a 2-byte
vector size and a 4-byte count. The format string \code{"<4sHHI"} describes exactly that.

\begin{center}
\begin{tikzpicture}[>=Stealth,decoration={brace,amplitude=3pt}]
  \foreach \h [count=\i from 0] in {45,4d,42,31,01,00,04,00,03,00,00,00} {
    \ifnum\i<4 \node[c18b=cblue] (x\i) at (\i*0.8,0) {\h};
    \else\ifnum\i<6 \node[c18b=cgreen] (x\i) at (\i*0.8,0) {\h};
    \else\ifnum\i<8 \node[c18b=cgold] (x\i) at (\i*0.8,0) {\h};
    \else \node[c18b=cpurple] (x\i) at (\i*0.8,0) {\h};
    \fi\fi\fi
    \node[font=\sffamily\tiny,text=cgray] at (\i*0.8,-0.55) {\i};}
  \draw[decorate,cblue,thick] (-0.4,0.45) -- node[above=2pt,font=\sffamily\scriptsize,text=ink]{magic: \code{4s}} (2.8,0.45);
  \draw[decorate,cgreen,thick] (3.2,0.45) -- node[above=2pt,font=\sffamily\scriptsize,text=ink]{version: \code{H}} (4.4,0.45);
  \draw[decorate,cgold,thick] (4.8,0.45) -- node[above=2pt,font=\sffamily\scriptsize,text=ink]{dim: \code{H}} (6.0,0.45);
  \draw[decorate,cpurple,thick] (6.4,0.45) -- node[above=2pt,font=\sffamily\scriptsize,text=ink]{count: \code{I}} (9.2,0.45);
  \node[font=\sffamily\scriptsize,text=cgray,anchor=west] at (9.5,0) {byte offsets below};
\end{tikzpicture}
\end{center}
% MD-BEGIN
% ```text
% offset :   0  1  2  3 | 4  5 | 6  7 | 8  9 10 11
% bytes  :  45 4d 42 31 |01 00 |04 00 |03 00 00 00      (hex)
% field  :  magic "EMB1"|version|  dim | count
% format :     4s      |   H   |  H   |    I           -> "<4sHHI"
% ```
% MD-END

\begin{pycode}
import struct

header = struct.pack("<4sHHI", b"EMB1", 1, 4, 3)
print(header)
print(len(header), struct.calcsize("<4sHHI"))
magic, version, dim, count = struct.unpack("<4sHHI", header)
print(magic, version, dim, count)
print(type(header), type(struct.unpack("<4sHHI", header)))
\end{pycode}
\begin{rawcode}
b'EMB1\x01\x00\x04\x00\x03\x00\x00\x00'
12 12
b'EMB1' 1 4 3
<class 'bytes'> <class 'tuple'>
\end{rawcode}

\begin{typecheck}
\code{pack} returns \code{bytes}. \code{unpack} always returns a \emph{tuple}, even when there is
one value, so write \code{(n,) = struct.unpack("<I", b)} or use \code{[0]}. A \code{4s} field comes
back as \code{bytes}, not \code{str}. Decode it yourself if it holds text.
\end{typecheck}

\subsection{The format string}

The \emph{first} character sets byte order (and whether sizes are standard). Then one letter per
field, optionally with a count in front: \code{"4f"} means four floats, \code{"2H"} means two
unsigned shorts.

\begin{center}
\begin{tabularx}{\linewidth}{@{}l l X@{}}
\toprule
\textbf{Char} & \textbf{Meaning} & \textbf{Python type} \\
\midrule
\code{<} \ \code{>} & little / big endian, standard sizes & (prefix) \\
\code{=} \ \code{@} & native order; \code{@} also adds alignment padding & (prefix) \\
\code{b} \ \code{B} & signed / unsigned byte (1) & \code{int} \\
\code{h} \ \code{H} & signed / unsigned short (2) & \code{int} \\
\code{i} \ \code{I} & signed / unsigned int (4) & \code{int} \\
\code{q} \ \code{Q} & signed / unsigned long long (8) & \code{int} \\
\code{f} \ \code{d} & float32 (4) / float64 (8) & \code{float} \\
\code{?} & bool (1) & \code{bool} \\
\code{Ns} & \code{N} raw bytes, padded with zero bytes & \code{bytes} \\
\code{x} & one padding byte, skipped & (no value) \\
\bottomrule
\end{tabularx}
\end{center}

\begin{gamerule}
Always start the format with \code{<} or \code{>}. Without a prefix Python uses your machine's
native order \emph{and} may insert padding bytes between fields, so \code{@bI} is 8 bytes while
\code{<bI} is 5. A file format is defined by its spec, not by your CPU.
\end{gamerule}

\subsection{Endianness}

A number bigger than one byte must be stored in some order. \emph{Big endian} puts the most
significant byte first, \emph{little endian} the least significant first. Most CPUs (x86, ARM
in normal mode) are little endian. Network protocols and many file headers are big endian.

\begin{center}
\begin{tikzpicture}[>=Stealth]
  \node[font=\sffamily\footnotesize,anchor=east] at (-0.5,1.5) {\code{0x12345678}, big endian};
  \node[font=\sffamily\footnotesize,anchor=east] at (-0.5,0) {little endian};
  \foreach \h [count=\i from 0] in {12,34,56,78} \node[c18b=cblue,minimum width=11mm] at (\i*1.2,1.5) {\h};
  \foreach \h [count=\i from 0] in {78,56,34,12} \node[c18b=cgold,minimum width=11mm] at (\i*1.2,0) {\h};
  \foreach \i in {0,1,2,3} \node[font=\sffamily\tiny,text=cgray] at (\i*1.2,0.65) {addr \i};
  \node[font=\sffamily\scriptsize,text=cgray,anchor=west] at (4.4,1.5) {most significant byte first};
  \node[font=\sffamily\scriptsize,text=cgray,anchor=west] at (4.4,0) {least significant byte first};
\end{tikzpicture}
\end{center}
% MD-BEGIN
% ```text
% value 0x12345678
% address        :   0    1    2    3
% big endian ">I":  12   34   56   78    (most significant byte first)
% little endian "<I": 78   56   34   12    (least significant byte first)
% ```
% MD-END

\begin{pycode}
print(struct.pack(">I", 0x12345678).hex(" "))
print(struct.pack("<I", 0x12345678).hex(" "))
print((0x12345678).to_bytes(4, "big").hex(" "))
print(int.from_bytes(b"\x01\x00", "little"), int.from_bytes(b"\x01\x00", "big"))
print(int.from_bytes(b"\xff\xff", "big", signed=True))
\end{pycode}
\begin{rawcode}
12 34 56 78
78 56 34 12
12 34 56 78
1 256
-1
\end{rawcode}

For a single integer, \code{int.to\_bytes()} and \code{int.from\_bytes()} are simpler than
\code{struct}. The same two bytes \code{01 00} are 1 in little endian and 256 in big endian, so a
wrong guess gives a wrong number, not an error.

\subsection{Reading from the middle of a buffer}

You rarely have just the header. \code{unpack\_from(fmt, buf, offset)} reads at an offset without
slicing. \code{iter\_unpack} walks over a repeated record. The \code{x} code skips bytes you do not
care about.

\begin{pycode}
vec = struct.pack("<4f", 0.5, -1.25, 3.0, 0.1)
blob = header + vec
print(struct.unpack_from("<4f", blob, 12))
print(struct.unpack_from("<4sHHI", blob))
for (x,) in struct.iter_unpack("<f", vec):
    print(x)
print(struct.unpack("<8xHH2x", b"12345678\x07\x00\x08\x00.."))
\end{pycode}
\begin{rawcode}
(0.5, -1.25, 3.0, 0.10000000149011612)
(b'EMB1', 1, 4, 3)
0.5
-1.25
3.0
0.10000000149011612
(7, 8)
\end{rawcode}

Look at the last float: \code{0.1} stored as float32 comes back as
\code{0.10000000149011612}. Float32 holds about seven decimal digits. This is why model weights
saved as float32 and read as Python \code{float} (which is float64) show these long tails. It is
not a bug.

\subsection{Typical errors}

\begin{pycode}
for bad in (lambda: struct.unpack("<I", b"\x01\x02"),
            lambda: struct.pack("<B", 300),
            lambda: struct.pack("<4s", "abcd")):
    try:
        bad()
    except struct.error as e:
        print(e)
print(struct.pack("<8s", b"abc"))
print(struct.calcsize("@bI"), struct.calcsize("<bI"))
\end{pycode}
\begin{rawcode}
unpack requires a buffer of 4 bytes
'B' format requires 0 <= number <= 255
argument for 's' must be a bytes object
b'abc\x00\x00\x00\x00\x00'
8 5
\end{rawcode}

\begin{trap}
Three classic mistakes: giving \code{unpack} a buffer of the wrong length (it must match
\code{calcsize} exactly), passing a \code{str} to an \code{s} field (you need \code{bytes}), and
packing a number that does not fit the type. All raise \code{struct.error}.
\end{trap}

\subsection{A real-world shape: length, then JSON, then data}

Many ML file formats start with a small fixed header that tells you how long the next, flexible
part is. The pattern below (a length, a JSON block, then raw floats) imitates this idea.

\begin{pycode}
import json

meta = json.dumps({"w": {"dtype": "F32", "shape": [2]}}).encode("utf-8")
blob = struct.pack("<Q", len(meta)) + meta + struct.pack("<2f", 1.5, -2.0)
(n,) = struct.unpack_from("<Q", blob)
info = json.loads(blob[8:8 + n])
print(n, info)
print(struct.unpack_from("<2f", blob, 8 + n))
print(type(info))
\end{pycode}
\begin{rawcode}
37 {'w': {'dtype': 'F32', 'shape': [2]}}
(1.5, -2.0)
<class 'dict'>
\end{rawcode}

\begin{boss}[: nothing in the bytes tells you the layout]
Binary data has no field names, no types, no separators. The same 4 bytes can be an
\code{int}, a \code{float}, four tiny numbers or text. The meaning lives in a spec, and your
format string is your copy of that spec. To decode safely: (1) check a magic value first,
(2) pick byte order explicitly, (3) pick the exact sizes and signedness, (4) check the total
length with \code{calcsize}. When a number comes out absurd (millions where you expect 3), suspect
byte order first.
\end{boss}

% ------------------------------------------------------------
\section{Extraction with Binary Data Tools}

\code{struct} gets clumsy when formats are nested or conditional. Third-party packages describe
the layout declaratively instead: \code{construct}, \code{bitstring}, \code{hachoir} and Kaitai
Struct are the ones to know. You describe the fields once and the library parses (and sometimes
builds) the bytes. Below is the same header with \code{construct}. It needs
\code{pip install construct}; I did not run it here, so check the library docs for the exact API
of your version.

\begin{pycode}
# needs: pip install construct   (not run in these notes)
from construct import Struct, Const, Int16ul, Int32ul

Header = Struct(
    "magic" / Const(b"EMB1"),
    "version" / Int16ul,      # ul = unsigned, little endian
    "dim" / Int16ul,
    "count" / Int32ul,
)
h = Header.parse(header)         # header = the 12 bytes built above
print(h.dim, h.count)
\end{pycode}

\begin{center}
\begin{tabularx}{\linewidth}{@{}l X X@{}}
\toprule
 & \textbf{\code{struct}} & \textbf{declarative library} \\
\midrule
Install & built in & \code{pip install} \\
Style & format string, positions by hand & named fields, nested records \\
Good for & flat, fixed records & complex or conditional formats \\
Readability & terse, easy to mis-type & self-documenting \\
\bottomrule
\end{tabularx}
\end{center}

% ------------------------------------------------------------
\section{Convert Bytes/Strings with binascii}

Raw bytes print as a mix of letters and \code{\textbackslash x} escapes. To \emph{see} them, convert
to hex. To \emph{send} them through a text-only channel such as JSON, convert to base64.

\begin{pycode}
import binascii, base64

print(binascii.hexlify(b"\x89PNG"))
print(binascii.unhexlify(b"89504e47"))
print(b"\x89PNG".hex(), bytes.fromhex("89 50 4e 47"))
print(type(b"ab".hex()), type(binascii.hexlify(b"ab")))
\end{pycode}
\begin{rawcode}
b'89504e47'
b'\x89PNG'
89504e47 b'\x89PNG'
<class 'str'> <class 'bytes'>
\end{rawcode}

\begin{typecheck}
\code{bytes.hex()} returns a \code{str}, and \code{bytes.fromhex(str)} takes one. The older
\code{binascii.hexlify} returns \code{bytes}. The methods on \code{bytes} are usually nicer. Pass
\code{hex(" ")} to put a space between bytes.
\end{typecheck}

Base64 turns any bytes into plain ASCII letters, digits, \code{+}, \code{/} and \code{=}. APIs that
accept an image inside JSON expect exactly this.

\begin{pycode}
raw = bytes(range(250, 256)) + b"hi"
b64 = base64.b64encode(raw)
print(b64, type(b64))
print(b64.decode("ascii"))
print(base64.b64decode(b64) == raw)
print(hex(binascii.crc32(b"hello")))
\end{pycode}
\begin{rawcode}
b'+vv8/f7/aGk=' <class 'bytes'>
+vv8/f7/aGk=
True
0x3610a686
\end{rawcode}

\begin{center}
\begin{tikzpicture}[box/.style={draw=#1,fill=#1!10,rounded corners=3pt,minimum height=8mm,
    font=\ttfamily\small,inner sep=4pt},>=Stealth]
  \node[box=cblue] (a) at (0,0) {bytes};
  \node[box=cgray] (b) at (4.2,0) {base64 bytes};
  \node[box=cgreen] (c) at (8.2,0) {str};
  \draw[->,thick] (a) -- node[above,font=\ttfamily\scriptsize]{b64encode} (b);
  \draw[->,thick] (b) -- node[above,font=\ttfamily\scriptsize]{.decode("ascii")} (c);
  \node[font=\sffamily\scriptsize,text=cgray] at (8.2,-0.8) {safe to put in JSON};
\end{tikzpicture}
\end{center}
% MD-BEGIN
% ```mermaid
% flowchart LR
%     A["bytes (image)"] -- "base64.b64encode" --> B["base64 bytes"]
%     B -- ".decode('ascii')" --> C["str (safe to put in JSON)"]
% ```
% MD-END

\begin{trap}
\code{b64encode} returns \code{bytes}, and \code{json.dumps} cannot serialize \code{bytes}. Call
\code{.decode("ascii")} to get a \code{str} first. Base64 is an \emph{encoding}, not encryption,
and it makes the data about one third larger.
\end{trap}

% ------------------------------------------------------------
\section{Use Bit Operators}

Integers are stored as bits. Python lets you work on those bits with six operators. They work on
\code{int} only.

\begin{center}
\begin{tabularx}{\linewidth}{@{}l l X l@{}}
\toprule
\textbf{Operator} & \textbf{Name} & \textbf{Result bit is 1 when ...} & \textbf{Example} \\
\midrule
\code{\&} & and & both bits are 1 & \code{5 \& 3 = 1} \\
\code{|} & or & at least one bit is 1 & \code{5 | 3 = 7} \\
\code{\^{}} & xor & exactly one bit is 1 & \code{5 \^{} 3 = 6} \\
\code{\~{}} & not & flips every bit (\code{\~{}x} is \code{-x-1}) & \code{\~{}5 = -6} \\
\code{<{}<} & left shift & bits move left, zeros fill in (times 2) & \code{5 <{}< 1 = 10} \\
\code{>{}>} & right shift & bits move right (floor-divide by 2) & \code{5 >{}> 1 = 2} \\
\bottomrule
\end{tabularx}
\end{center}

\begin{center}
\begin{tikzpicture}[>=Stealth]
  \foreach \lab/\y/\col/\ba/\bb/\bc/\bd in {x = 5/0/cblue/0/1/0/1, y = 3/-0.8/cgreen/0/0/1/1,
        x \& y = 1/-1.8/cred/0/0/0/1, x | y = 7/-2.6/cred/0/1/1/1, x \^{} y = 6/-3.4/cred/0/1/1/0} {
    \node[font=\ttfamily\small,anchor=east] at (-0.4,\y) {\lab};
    \node[c18bit,fill=\col!12,draw=\col] at (0,\y) {\ba};
    \node[c18bit,fill=\col!12,draw=\col] at (0.7,\y) {\bb};
    \node[c18bit,fill=\col!12,draw=\col] at (1.4,\y) {\bc};
    \node[c18bit,fill=\col!12,draw=\col] at (2.1,\y) {\bd};
  }
  \draw[cgray,dashed] (-0.4,-1.3) -- (2.5,-1.3);
\end{tikzpicture}
\end{center}
% MD-BEGIN
% ```text
% x      = 5   0 1 0 1
% y      = 3   0 0 1 1
%              -------
% x & y  = 1   0 0 0 1     1 only where both are 1
% x | y  = 7   0 1 1 1     1 where at least one is 1
% x ^ y  = 6   0 1 1 0     1 where exactly one is 1
% ```
% MD-END

\begin{pycode}
x, y = 0b0101, 0b0011
print(x & y, x | y, x ^ y, ~x, x << 1, x >> 1)
print(bin(x), format(x, "08b"), f"{x:#06b}", f"{255:#x}")
print(~x, ~x & 0xF, f"{~x & 0xFF:08b}")
print(-1 & 0xFF, -6 & 0xFF)
\end{pycode}
\begin{rawcode}
1 7 6 -6 10 2
0b101 00000101 0b0101 0xff
-6 10 11111010
255 250
\end{rawcode}

\begin{typecheck}
Python's \code{int} has no fixed width, so \code{\~{}5} is \code{-6}, not a big positive number.
Use a mask to see a fixed width: \code{\~{}5 \& 0xFF} is \code{250}. Negative numbers behave as if
they had infinitely many leading 1 bits (two's complement), which is why \code{-6 \& 0xFF} gives
250.
\end{typecheck}

\subsection{Using bits as flags and fields}

A common use is packing several yes/no options into one integer, or pulling a field out of a
byte with a shift and a mask.

\begin{pycode}
READ, WRITE, EXEC = 0b100, 0b010, 0b001
perm = READ | WRITE              # set two flags
print(perm, bool(perm & WRITE), bool(perm & EXEC))
perm |= EXEC                     # set a flag
perm &= ~WRITE                   # clear a flag
print(f"{perm:03b}")
perm ^= READ                     # toggle a flag
print(f"{perm:03b}")
b = 0b10110100
print((b >> 4) & 0xF, b & 0xF)   # high nibble, low nibble
print((255).bit_count(), (0b1011).bit_length())
\end{pycode}
\begin{rawcode}
6 True False
101
001
11 4
8 4
\end{rawcode}

\begin{center}
\begin{tikzpicture}[>=Stealth]
  \foreach \b [count=\k from 0] in {1,0,1,1}
    \node[c18bit,fill=cblue!12,draw=cblue] at (\k*0.7,0) {\b};
  \foreach \b [count=\k from 4] in {0,1,0,0}
    \node[c18bit,fill=cgold!12,draw=cgold] at (\k*0.7,0) {\b};
  \draw[decorate,decoration={brace,mirror,amplitude=3pt},cblue,thick] (-0.3,-0.45) -- node[below=2pt,font=\sffamily\scriptsize,text=ink]{\code{b >{}> 4} = 11} (2.5,-0.45);
  \draw[decorate,decoration={brace,mirror,amplitude=3pt},cgold,thick] (2.6,-0.45) -- node[below=2pt,font=\sffamily\scriptsize,text=ink]{\code{b \& 0xF} = 4} (5.4,-0.45);
\end{tikzpicture}
\end{center}
% MD-BEGIN
% ```text
% b = 1 0 1 1 | 0 1 0 0
%     high nibble  low nibble
%     b >> 4 = 1011 = 11      b & 0xF = 0100 = 4
% ```
% MD-END

\begin{trap}
Do not confuse \code{\&} with \code{and}. \code{6 and 3} is \code{3} (logical: the last truthy
value), \code{6 \& 3} is \code{2} (bit by bit). And \code{<{}<} with a huge shift quietly builds a
huge integer. It never overflows.
\end{trap}

% ------------------------------------------------------------
\section{Review and what comes next}

You can now describe a binary layout with a format string, read and write it with \code{struct},
display bytes as hex, move them through JSON as base64, and flip single bits. The next chapter
turns to dates and times.

\begin{practice}
\begin{enumerate}[nosep,leftmargin=*]
  \item Write the hex string \code{"4749463839"} to a \code{bytes} variable with
        \code{bytes.fromhex}. Does it start with the ASCII letters \code{GIF}? Show how you check.
  \item Pack the three integers 7, 300 and 70000 as unsigned 1-, 2- and 4-byte little-endian
        values. What is the length of the result, and which format string did you use?
  \item The bytes \code{b"\textbackslash x00\textbackslash x00\textbackslash x01\textbackslash x00"}
        are an unsigned 4-byte integer. What are its values read big endian and little endian?
  \item Why does \code{struct.unpack("<I", data)} raise an error when \code{data} has 5 bytes?
        Show two ways to read only the first four.
  \item Store the flags \code{ADMIN=1}, \code{BETA=2}, \code{DARK=4} in one integer. Turn on
        \code{BETA}, turn off \code{ADMIN}, and test for \code{DARK}.
  \item Take 10 random bytes, convert to base64, put it in a dict, dump it with \code{json}, load
        it back and recover the original bytes.
\end{enumerate}
\end{practice}

\begin{achievement}{Byte Wrangler}
\begin{itemize}[nosep,leftmargin=*]
  \item Read and write fixed binary layouts with \code{struct}, choosing byte order on purpose.
  \item Explain endianness and spot it as the cause of absurd numbers.
  \item Show bytes as hex, and ship them inside JSON as base64.
  \item Set, clear, toggle and test bits, and extract a bit field with shift and mask.
\end{itemize}
\end{achievement}
